\subsection{连通交换实李群}

令 \(G\) 为连通交换（实）李群。指数映射

\[
\exp_ {\mathrm{G}}: \mathrm{L} (\mathrm{G}) \to \mathrm{G}
\]

是李群态射，且是满射、有离散核（第III章§6第4节命题11），由此可知 \(\mathrm{L}(\mathrm{G})\) 是 \(\mathrm{G}\) 的连通覆盖。

\begin{enumerate}
\def\labelenumi{\alph{enumi})}
\item
  下列条件等价：G 单连通，\(\exp_{G}\) 是同构，G 同构于 \(R^{n}\)（\(n = \dim G\)）。在这种情形，用同构 \(\exp_{G}\) 把 L(G) 的向量空间结构转移到 G 上，就得到 G 上的一个向量空间结构，它是唯一与 G 的拓扑群结构相容的向量空间结构。单连通交换李群称为向量（李）群；除另有说明外，它们总被赋予上面定义的 R-向量空间结构。
\item
  以 \(\Gamma(\mathrm{G})\) 表示 \(\exp_{G}\) 的核。由《一般拓扑学》第VII章§1第1节定理1，群 G 紧当且仅当 \(\Gamma(\mathrm{G})\) 是 \(\mathrm{L}(\mathrm{G})\) 中的格，换句话说（同处所引），当且仅当自由 Z-模 \(\Gamma(\mathrm{G})\) 的秩等于 G 的维数。反之，若 L 是有限维 R-向量空间，\(\Gamma\) 是 L 中的格，则商拓扑群 \(L/\Gamma\) 是紧连通交换李群。
\end{enumerate}

紧连通交换李群称为实环面，或（在本章中）环面。

\begin{enumerate}
\def\labelenumi{\alph{enumi})}
\setcounter{enumi}{2}
\tightlist
\item
  在一般情形，令 E 为 L(G) 中由 \(\Gamma(\mathrm{G})\) 生成的向量子空间，令 V 为其补子空间。则 G 是它的李子群 \(\exp(\mathrm{E})\) 与 \(\exp(\mathrm{V})\) 的直积；前者是环面，后者是向量群。最后，G 的每个紧子群都含于 \(\exp(\mathrm{E})\)（因为它在 \(\exp(V)\) 上的投影必然化为单位元）；于是，子群 \(\exp(E)\) 是 G 的唯一的极大紧子群。
\end{enumerate}

例如，取 \(\mathrm{G} = \mathrm{C}^*\)；把 \(\mathrm{L(G)}\) 与 \(\mathbf{C}\) 等同起来，使 \(\mathrm{G}\) 的指数映射为 \(x \mapsto e^x\)。则 \(\Gamma(\mathrm{G}) = 2\pi i\mathbf{Z}, \mathrm{E} = i\mathbf{R}\)，从而 \(\exp(\mathrm{E}) = \mathbf{U}\)；若取 \(\mathrm{V} = \mathbf{R}\)，则 \(\exp(\mathrm{V}) = \mathbf{R}_+^*\)，我们重新得到《一般拓扑学》第VIII章§1第3节命题1 中构造的同构 \(\mathbf{C}^* \to \mathbf{U} \times \mathbf{R}_+^*\)。

\(d\) ) 最后注意，\(\exp_{\mathrm{G}}:\mathrm{L}(\mathrm{G})\to \mathrm{G}\) 是 \(\mathbf{G}\) 的泛覆盖，因而 \(\Gamma (\mathbf{G})\) 可自然地与 \(\mathbf{G}\) 的基本群等同起来。

